\pendiente{Ajustar al formato que exija el ED-TFG para la declaración de uso de IA
(plataforma de TFG, art.~5.1.3 de la normativa \cite{eps_normativa_tfg_2026}); el contenido
de este apéndice está pensado para poder trasladarse a ese formato sin cambios de fondo.}

Este apéndice detalla la declaración resumida en la \cref{SEC:IMP-AI}. Todo lo que afirma
puede comprobarse en el repositorio del \dfn{frontend} (historial de Git, ficheros de
instrucciones y registros de coordinación) o, cuando se trata de la orquestación, en los
\textit{scripts} del autor que se citan, que residen fuera del repositorio. Las cifras se
refieren al estado de las ramas a 19 de octubre de 2026 (\textit{commits} \texttt{a21f8ef}
de la línea del núcleo y \texttt{91475a0} de la línea de gestión).

\section*{Herramientas}

\begin{description}
  \item[Claude Code (Anthropic).] Agente de programación de línea de órdenes
        \cite{claude_code_docs}, ejecutado en la máquina del autor sobre el
        \dfn{worktree} \texttt{SCDE-claude} (rama \texttt{claude/core-flows}). En la fase
        final se usó con el modelo Claude Opus~5.5, que es el que figura en las líneas de
        coautoría de los \textit{commits} descritas más abajo. Ya se había utilizado en el
        prototipo previo: el \textit{commit} \texttt{1e8e9bd} (15 de abril de 2026) añade
        ocho subagentes en \texttt{.claude/agents/} (análisis del \dfn{backend}, capa de
        servicios, autenticación, anotaciones, paneles, revisión de código y redacción),
        que siguen en el repositorio aunque el prototipo no forma parte del resultado.
  \item[Codex (OpenAI).] Agente de programación \cite{openai_codex}, configurado con el
        modelo \texttt{gpt-6.1-sol}, sobre el \dfn{worktree} \texttt{SCDE-codex} (rama
        \texttt{codex/admin-ops}). Se ejecutaba en un entorno aislado que no le permitía
        escribir en el directorio \texttt{.git}; por eso, en la fase autónoma, redactaba
        el mensaje de cada \textit{commit} en un fichero y el lanzador lo consolidaba al
        terminar, limitado a las rutas que el mapa de propiedad asigna a su línea.
\end{description}

Ninguna de las dos herramientas tuvo acceso a credenciales de publicación en repositorios
remotos: las reglas del proyecto prohíben a ambas hacer \textit{push}, reescribir el
historial o fusionar ramas ajenas, y el único \textit{merge} permitido es el de la rama local
de integración, que solo combina las dos líneas para ejecutar la batería conjunta.

\section*{Instrucciones de proyecto}

Los agentes no recibieron indicaciones sueltas sino un conjunto de ficheros versionados que
leían al comenzar cada sesión (\cref{TAB:AI-INSTRUCTIONS}). Al estar en el repositorio,
cualquier cambio de reglas queda registrado como un \textit{commit} más y sujeto al mapa de
propiedad.

\begin{table}[htbp]
  \centering
  \small
  \begin{tabular}{p{0.34\textwidth}p{0.08\textwidth}p{0.48\textwidth}}
    \toprule
    \textbf{Fichero} & \textbf{Líneas} & \textbf{Función} \\
    \midrule
    \texttt{AGENTS.md} & 16 & Reglas comunes: lecturas previas, rutas propias, no inventar
      operaciones, simulación con la forma exacta del OpenAPI, verificación antes del
      \textit{commit}, prohibiciones de Git y mensajes sin mención de herramientas de
      \acs{ia}. \\
    \texttt{CLAUDE.md} & 21 & Contexto y convenciones específicos de Claude Code. \\
    \texttt{docs/coordination/CODEX\_BRIEF.md} & 98 & Encargo priorizado de la línea de
      gestión: alcance por módulos, orden de lectura y criterios de aceptación. \\
    \texttt{docs/coordination/OWNERSHIP.json} & --- & Mapa de propiedad de rutas por
      línea, que comprueba \texttt{ownership-check.mjs}. \\
    \texttt{docs/coordination/REQUESTS.md} & --- & Canal de peticiones entre líneas para
      cambios en ficheros compartidos y hallazgos de auditoría. \\
    \texttt{docs/frontend/STYLE\_GUIDE.md} & 73 & Lenguaje visual v2, de obligado
      cumplimiento (\cref{AP:STYLE}). \\
    \texttt{.claude/skills/scde-frontend/SKILL.md} & 54 & Procedimiento para construir un
      módulo y autocrítica visual \cite{anthropic_skills_frontend_design}. \\
    \texttt{.claude/agents/*.md} & --- & Nueve subagentes acotados; ocho del prototipo y un
      auditor de solo lectura de la rama paralela, añadido en esta fase. \\
    \bottomrule
  \end{tabular}
  \caption{Ficheros de instrucciones que leían los agentes (recuento de líneas en
    \texttt{a21f8ef}).\label{TAB:AI-INSTRUCTIONS}}
\end{table}

En la fase de iteraciones autónomas se añadió un encargo de misión, fuera del repositorio,
con la lista priorizada de tareas de cada agente, las reglas de cierre de una iteración y un
fichero de progreso por agente, en el que cada ejecución anotaba qué había hecho, con qué
\textit{commit} y qué seguía. Ese encargo remite a las reglas del repositorio y declara que,
en caso de conflicto, prevalecen estas.

\section*{Orquestación del trabajo autónomo}

El autor lanzaba y seguía las iteraciones desde un servicio de mensajería: un programa en
un ordenador de placa reducida recibía las órdenes y, mediante una conexión SSH restringida
a una única orden, invocaba en la máquina de desarrollo un lanzador que arrancaba la sesión
del agente sobre su \dfn{worktree}, guardaba su registro y, al terminar, la relanzaba con la
tarea siguiente. Cada agente enviaba un mensaje breve al empezar y otro al acabar, con el
\textit{commit} resultante, y el autor podía detener una ejecución en cualquier momento.

El diseño inicial de este relanzador permitió que se ejecutaran varias sesiones del mismo
agente a la vez sobre el mismo \dfn{worktree} (\cref{TAB:IT-PHASE2}, caso~31). Lo detectó
el autor al revisar el historial y se corrigió en el propio lanzador, que desde entonces
rechaza una segunda sesión de un agente mientras la anterior siga activa. La lección se
recoge en la \cref{SEC:GEN-PROPOSAL}.

\section*{Reparto del trabajo}

La \cref{TAB:AI-SPLIT} resume quién produjo cada parte. En la columna de revisión,
\emph{cruzada} significa que la otra línea auditó el resultado con las comprobaciones de la
\cref{SEC:TEST-AUDITS}; la revisión final del autor sobre todo el contenido está pendiente
de su propia declaración al final de la \cref{SEC:IMP-AI}.

\begin{table}[htbp]
  \centering
  \small
  \begin{tabular}{p{0.30\textwidth}p{0.26\textwidth}p{0.34\textwidth}}
    \toprule
    \textbf{Parte} & \textbf{Producida por} & \textbf{Verificación} \\
    \midrule
    Análisis del contrato e inventario de la \acs{api} & Claude Code (\textit{scripts}
      deterministas) & Contraste con las vistas de Django en los casos dudosos \\
    Base común y sistema de diseño & Claude Code & Batería, auditor de estilo; dirección
      visual aceptada por el autor tras rechazar la primera \\
    Núcleo docente (corrección, exámenes, escaneos, resultados, revisiones) & Claude Code
      (34 \textit{commits} desde \texttt{1ebeac9}) & Batería propia y conjunta, prueba de
      contrato \\
    Gestión (organización, asignaturas, cuenta, notificaciones) & Codex (28
      \textit{commits}) & Auditoría cruzada por Claude Code, prueba de contrato en la rama
      de integración \\
    Publicación de la demostración & Codex, con autorización del autor & Comprobación de
      la dirección publicada antes de anunciarla \\
    Capítulo 2 (estado del arte) & Borradores de Codex & Revisión e integración por
      Claude Code \\
    Resto de la memoria y apéndices & Claude Code; partes de gestión desde borradores de
      Codex & Contraste con el código, el historial y \texttt{docs/tfg/} \\
    \bottomrule
  \end{tabular}
  \caption{Reparto del trabajo entre los agentes.\label{TAB:AI-SPLIT}}
\end{table}

\section*{Decisiones del autor}

Las decisiones que no tomaron los agentes, por estar reservadas al autor en las
instrucciones o por requerir su criterio, fueron: el alcance y los objetivos del trabajo
(\cref{CAP:INTRODUCTION}); la elección de las dos herramientas y de la organización en dos
líneas paralelas con un auditor; las reglas de Git y de propiedad de ficheros; el rechazo
de la primera dirección visual (\cref{TAB:IT-PHASE1}, caso~16); la autorización para
publicar la demostración; y el diseño y la corrección del relanzador. Las decisiones menores
que surgieron durante las iteraciones autónomas las tomaron los agentes con un criterio
explícito y quedaron anotadas en los registros (\texttt{docs/tfg/}, los ficheros de progreso
y el registro de la línea de gestión) para su revisión posterior.

\section*{Errores introducidos por los agentes}

Los registros recogen los siguientes errores atribuibles a los agentes, con el mecanismo que
los detectó; el detalle de cada uno está en el \cref{AP:ITERATIONS}:
\begin{itemize}
  \item \emph{De proceso}: un \textit{commit} con pruebas en rojo por no comprobar el código
        de salida (caso~13); una afirmación documentada sobre la auditoría de propiedad que
        no se correspondía con el código (caso~24); y un \textit{commit} de la línea de
        gestión que importaba un fichero no incluido y rompió la rama de integración
        (caso~27). Los dos primeros los detectaron la propia batería y el contraste de la
        memoria con el código; el tercero, la batería de la rama de integración.
  \item \emph{De diseño}: una dirección visual que imitaba la paleta de la referencia en
        lugar de su principio (caso~16), detectada por el autor.
  \item \emph{De funcionamiento}: los dos defectos del autoguardado (casos 21 y 22), la
        pérdida de respuestas al recargar la bandeja de revisiones (caso~29), las tablas
        que se desmontaban al refrescar (caso~30) y una regresión del editor de permisos de
        la línea de gestión. Todos se corrigieron con una prueba de regresión.
  \item \emph{De las herramientas propias}: \texttt{coverage.py} infravaloraba la cobertura
        (caso~23) y el recorte del OpenAPI generaba falsos positivos (caso~25).
  \item \emph{De documentación}: un recuento de operaciones del contrato (96 en lugar de
        95) y un párrafo de un borrador del estado del arte que describía carencias ya
        corregidas, ambos detectados al revisar la memoria.
\end{itemize}

Un último error afecta a esta propia declaración. La regla de que los mensajes de
\textit{commit} no mencionen herramientas de \acs{ia} no se cumplió en 14
\textit{commits} de la línea del núcleo, que incluyen una línea
\texttt{Co-Authored-By} con el nombre del modelo, añadida por la configuración por defecto
de Claude Code. Se detectó al revisar el historial y desde entonces se omite; los
\textit{commits} afectados no se han modificado porque las reglas prohíben reescribir el
historial. La inconsistencia no oculta nada: esas líneas atribuyen el trabajo a la misma
herramienta que declara este apéndice.

\section*{Qué no se delegó en la verificación}

Ningún resultado se aceptó por la sola afirmación de un agente. Las cifras de pruebas y de
cobertura que aparecen en la memoria se obtuvieron ejecutando las órdenes del proyecto
sobre un \textit{commit} identificado, y cualquier lector puede reproducirlas con esas
mismas órdenes (\cref{SEC:IMP-TOOLS}). Lo que no se ha podido verificar se declara como
tal: no hubo pruebas con usuarios ni validación completa contra el \dfn{backend} real
(\cref{SEC:TEST-REAL}).
